
\problem{1}
Можно ли монетами в 10 и 15 сиклей набрать ровно 2019 сиклей?

\problem{2}
[CAMMAT, 2018, 9 класс] На доске написано три числа. За один ход разрешается стереть любые два числа $a$ и $b$ и вместо них написать числа $\frac{a}{b}$ и $b^2$. Можно ли с помощью таких операций из тройки $\left(\frac{1}{\sqrt{3}}, 1,1+\sqrt{3}\right)$ получить тройку $\left(\sqrt{3}, 3, \frac{1}{3+\sqrt{3}}\right)$ ?

\problem{3}
Гарри Поттер дрессирует шоколадную лягушку. Он заставляет прыгать ее вдоль числовой прямой. Первый прыжок лягушка совершает на 1 дециметр, следующий - на 2 дециметра, третий - на 3 дециметра, и так далее, 101-й прыжок - на 101 дециметр. При этом прыгать лягушка может в любую из двух сторон. Может ли в итоге лягушка оказаться в той же точке, в которой она начинала путешествие?

\problem{4}
[Миссия выполнима, 2020] На доске написаны все натуральные числа от 1 до 100. Можно любую пару чисел $x, y$ заменять на $x y-29 x-29 y+870$. Какое число останется после 99 таких операций?


\problem{5}
У Гермионы есть 64 шоколадные лягушки. При этом одна из них сделана из белого шоколада, а остальные 63 - из черного. Гермиона рассадила всех лягушек в клетки шахматной доски $8 \times 8$, по одной лягушке в каждую клетку. Выбирая строку или столбец доски и произнося заклинание "Шоколадус перекратус", Гермиона перекрашивает всех белых лягушек этой линии в черный цвет, а всех черных - в белый. Может ли Гермиона, произнося заклинание несколько раз, сделать все 64 лягушки черными?

\problem{6}
На доске вначале было записано $n$ чисел: $1,2, \ldots, n$. Разрешается стереть любые два числа на доске, а вместо них записать модуль их разности. Какое наименьшее число может оказаться на доске после $n-1$ таких операций (a) при $n=111$; (b) при $n=110$ ?